\subsection*{Part 3: Quick quiz, box or bust?}
Python accepts \texttt{lives}, \texttt{player\_1} and \texttt{bestTime}. \texttt{high score} has a space, \texttt{1st\_place} starts with a digit, and \texttt{import} is a word Python already uses.

\subsection*{Part 3, Mission: Score counter}
\py|score += 10| does exactly the same as \py|score = score + 10|, so nothing changes. To count down, start with \py|score = 50| and use \py|score -= 10|: the last line says \texttt{Final score: 0}. With \py|lives -= 1| inside the loop, the lives go 2, 1, 0, then $-1$ and $-2$. Python doesn't stop at 0 unless you tell it to (you will learn how with \py|if| in Part 5).

\textbf{Super challenge:}
\begin{codenonum}
import time
score = 0
points = 10
print("Score:", score)
for i in range(5):
    time.sleep(1)
    score = score + points
    points = points + 10
    print("Score:", score)
print("Game over! Final score:", score)
\end{codenonum}
The score goes 10, 30, 60, 100, 150. The final score is \texttt{150}.

\subsection*{Part 3, Mission: Growing spiral}
With \texttt{+ 8} the last line is 317 steps long, and the spiral only just fits on the stage. For the thicker pen, put \py|thick = 1| next to \py|size = 5|, and add \py|pensize(thick)| and \py|thick = thick + 0.2| inside the loop.

\textbf{Super challenge} (after the three starting lines):
\begin{codenonum}
speed(13)
pencolor("purple")
penup()
goto(-100, 100)
pendown()
size = 200
for i in range(40):
    forward(size)
    right(90)
    size = size - 5
\end{codenonum}

\subsection*{Part 3, Mission: Rainbow from a colour list}
\py|colours.append("pink")| adds pink to the end of the list, so the loop goes round 7 times and pink is the bottom stripe. Red, white and blue stripes, from top to bottom, make the flag of the Netherlands. Try \py|["white", "blue", "red"]| for Russia.

\textbf{Super challenge:} lines 8 to 16 become
\begin{codenonum}
y = 60
length = 300
for colour in colours:
    pencolor(colour)
    penup()
    goto(-150, y)
    pendown()
    forward(length)
    y = y - 15
    length = length - 30
hideturtle()
\end{codenonum}

\subsection*{Part 3, Mission: Random story maker}
For the places, add \py|places = ["in the bath", "on the bus"]| before the loop, \py|place = random.choice(places)| inside it, and change the print to \py|print(hero, action, thing, place)|. After the loop, \py|len(things)| is \texttt{4}, and \texttt{"my homework"} is the last item, because \py|append| always adds to the end:
\begin{outputbox}
4\\
{[}'a banana', 'the moon', 'a sock', 'my homework']
\end{outputbox}

\subsection*{Part 3: Translate it yourself}
\begin{codenonum}
pets = []
pets.append("cat")
pets.append("dog")
pets.append("fish")
print(len(pets))
print(pets[1])
\end{codenonum}
It prints \texttt{3} and then \texttt{dog}. Scratch's item 2 is Python's \texttt{[1]}. \py|pets = []| does the job of \chip[lists]{delete all of pets}, because it starts a brand-new empty list.

\subsection*{Part 3: Two-line puzzle}
\begin{codenonum}
for word in words:
    count += 1
\end{codenonum}
The loop goes round once for every item, and each time it adds 1 to \texttt{count}. That is exactly how \py|len| counts!
